\section{Funktionale Anforderungen}

% Fertigen Sie hier eine Liste aller funktionalen Anforderungen an. Halten Sie sich dabei an die Formatierung aus der Vorlesung und Übung. Sortieren Sie die Anforderungen nach Priorität. Die Priorität erhalten Sie aus dem Video. Falls etwas unklar ist, fragen Sie Ihren Customer-Proxy.

\subsection{Use Case-Diagramm}

% Viele der folgenden Unterkapitel sind leer. Füllen Sie sie entsprechend Ihren SWT-Kenntnissen; bedenken Sie dabei auch, dass ein UML-Diagramm oder auch ein Use Case ohne Erläuterung oft nicht aussagekräftig ist. Scheuen Sie sich also nicht, noch Kommentare oder weitere Angaben zu machen, wenn das der Klärung dient

\pagebreak

\subsection{Use Case-Beschreibungen}

% Abläufe mit interessanter Interaktion zwischen einem Akteur und verschiedenen Systemteilen werden als Use Cases beschrieben. Zu jedem Use Case kann es weitere, detaillierende Anforderungen geben, die dort angegeben sind (z.B R412). Rein interne Abläufe werden unter Umständen anders notiert (z.B. Ablauf- oder Sequenzdiagramm). Hier handelt es sich um Anforderungen des Kunden.


%%%% BEGIN - VORLAGE FÜR USE CASES %%%%
% Erstellen Sie eine eigene Datei für jeden Use Case

\subsubsection{UC: TITEL}

\begin{tabularx}{\linewidth}{|p{3.5cm}|X|}
	\hline
	\textbf{UC <01>}		& \textbf{<Titel>} \\ \hline
	\textbf{Umfeld}			& ... \\ \hline
	\textbf{Systemgrenze}	& ... \\ \hline
	\textbf{Ebene}			& ... \\ \hline
	\textbf{Hauptakteur}	& ... \\ \hline
	\textbf{Stakeholder}
        & Stakeholder1 möchte ...  \\
        & Stakeholder2 möchte ... \\
        & Stakeholder3 möchte ... \\
        \hline
	\textbf{Voraussetzungen}   & ... \\ \hline
	\textbf{Garantie}          & ... \\ \hline
	\textbf{Erfolgsfall}       & ... \\ \hline
	\textbf{Auslöser}          & ... \\ \hline	
	\textbf{Hauptszenario}			
        & 1. \\
	    & 2. \\
		& 3. \\
		& ... \\
        \hline
	\textbf{Erweiterung}			
        & 1a. WENN ... DANN ... \\ 
		& 2a. WENN ... DANN ... \\
        \hline
	\textbf{Priorität}     & ... \\ \hline
	\textbf{Technologie}   & ... \\ \hline
\end{tabularx}

\\[0.5cm]
\textbf{Erläuterungen und Details}
\begin{itemize}
    \item 
\end{itemize}

%%%% END - VORLAGE FÜR USE CASES %%%%